\section{Conclusion}
\label{sec:conclusion}

Run-level assurance is clearer when the claim is separated from the mechanism
and from capture resilience. G1--G4 identify anchored manifest integrity,
credential attribution, policy-relative coverage, and computation conformance;
V1--V4 are their cumulative prefixes. The cross-field map shows why a
signature, succinct proof, or attestation report contributes to particular
conditions without automatically constituting a complete level.

The executable study yields three lessons. Short authenticated records can be
cheaper than commitment entries. Sampling saves disclosure only under a stated
false-acceptance probability. In the demonstrated proof, credential attribution
dominates the scalar fusion arithmetic. These findings compare protocol
evidence, not deployability on a physically exposed device.

For deployment: authenticate and version $P$; choose the V-claim and tolerance;
state the capture profile; cover every required G-condition; and account for
online evidence, setup, retention, and judge work separately. Future work
should test an order-sensitive multi-sensor operator, a real attested path under
physical evaluation, and adaptive or dynamic-window adversaries.
